% =====================================================================================================
\section{Introduction}\label{sec:intro}
Home cameras, video doorbells and baby monitors send video over Wi-Fi and shared networks that their owners do not control. Two properties matter for such a link. \emph{Confidentiality} keeps a bystander on the network from watching the home; \emph{integrity} lets the viewer trust that what it shows was captured by the camera and was not edited, replayed or injected on the way. Encrypting the transport with TLS, SRTP~\cite{rfc3711} or a VPN provides both, but it hides the stream from every recorder, player and cloud service outside the key infrastructure, and the cipher work grows with the full bitrate of a low-power device.

Selective encryption encrypts only the parts of the compressed stream that carry most of the picture and keeps the format~\cite{massoudi2008}, so standard decoders still parse the result. For baseline JPEG, the sign and the appended amplitude bits of the quantized DCT coefficients can be encrypted without changing any Huffman code length~\cite{vandroogenbroeck2002,auer2013}. Such schemes are easy to weaken: the positions of the non-zero coefficients stay visible, so an attacker still sees the outline of the scene~\cite{massoudi2008,lin2024,hirose2024}; most use key material that does not change from frame to frame, so one known frame exposes the next; and almost none authenticate what they encrypt.

We ask how much of a secure camera link fits into the encoder hardware, and what an eavesdropper still learns. Our contributions are:
\begin{enumerate}
\item \textit{Authenticated selective encryption inside the entropy coder.} A Verilog MJPEG encoder XORs a keystream of the NIST standard Ascon-AEAD128~\cite{sp800232,ascon} onto coefficient signs and amplitudes. We prove that every code length is preserved (Lemma~\ref{lem:format}) and that the cipher never stalls the encoder (Proposition~\ref{prop:supply}); a 128-bit tag binds the scan to the protection level, session, frame counter and size.
\item \textit{A hardware form of a keyed macroblock permutation (E5).} The fabric codes one restart segment per macroblock with its own keystream and authenticates them in raster order; the ARM core sends them in a keyed order that changes with every frame. No frame buffer is needed, and the order is authenticated.
\item \textit{Parallel lanes and a measured system.} Restart segments make lanes independent; one, two and four lanes produce identical bytes. On a PYNQ-Z2 board every image is bit-exact with the golden model in every level, and the encoding clocks counted by the fabric equal the RTL simulation's to within \BdBusyVsRtlMaxPct\,\% (four lanes: \BdFourBusyMpix\,Mpixel/s at \BdFourMhz\,MHz).
\item \textit{One benchmark for all schemes.} Six protection levels and four published format-compliant schemes~\cite{vandroogenbroeck2002,kishore2015,auer2013,hirose2024} under the same keyless attacks, including puzzle solvers, sign recovery, a chosen-plaintext attack and trained networks, on \NumVideoClips{} clips and the Kodak set.
\item \textit{A closed feedback loop.} An authenticated feedback channel, a forward-secure key ratchet and an adaptive policy raise protection when the receiver reports an attack, in simulation and with the board as the camera.
\end{enumerate}
\TODO{Once the conference version is accepted: cite it here and state what this article adds (proofs, board evaluation, full benchmark, protocol analysis).}

% =====================================================================================================
\section{Related Work}\label{sec:related}
\textit{Selective encryption of JPEG.} Van Droogenbroeck and Benedett encrypted the appended bits of AC coefficients and kept the Huffman code words~\cite{vandroogenbroeck2002}; Puech and Rodrigues encrypted selected DCT coefficients with AES in OFB mode~\cite{puech2005}; Auer et al.\ permuted code-word/value pairs and toggled value bits at bitstream level at more than 25 VGA frames per second~\cite{auer2013}. Massoudi et al.\ note that visual distortion does not measure security~\cite{massoudi2008}. The idea has been carried to H.264/AVC and HEVC~\cite{boyadjis2017}.

\textit{Attacks.} Lin et al.\ recover DCT sign bits by linear programming~\cite{lin2024}. Hirose et al.\ show that the number of non-zero coefficients per block, which no field encryption hides, rebuilds contours unless the restart segments are permuted~\cite{hirose2024}; Li et al.\ recover block shuffles from one designed plaintext~\cite{li2025shuffle}; a GAN-based ciphertext-only attack recovers content from most privacy-preserving image schemes~\cite{sirichotedumrong2020}; and block permutations invite jigsaw solvers, from greedy placement~\cite{pomeranz2011} to minimum-spanning-tree assembly~\cite{gallagher2012}.

\textit{Hardware.} Kishore et al.\ added block and symbol scrambling to an FPGA JPEG encoder at 13--17\,\% more ALUTs~\cite{kishore2015}; Ba\u{g}baba and \"Ors combined an FPGA JPEG encoder with TEA encryption of DC and low AC coefficients~\cite{bagbaba2015}; Qiu et al.\ XOR raw frames with chaotic sequences on a satellite payload~\cite{qiu2025}. None authenticates frames, and their key material is fixed. Commercial JPEG cores process one sample per clock at 150--166\,MHz without cryptography~\cite{logijpge,castjpeg}. Ascon is small in hardware, about 7\,kGE for one unrolled round, and a threshold implementation costs about 3.1 times the area of the low-area core~\cite{gross2017}; Ascon-based MACs~\cite{khairallah2023} and Ascon with replay detection~\cite{kurian2025} have been built on Artix-7. Farahmand et al.\ measure cryptographic cores on a Zynq-7020 with DMA, cycle-exact timers and bit-exact comparison with software~\cite{farahmand2016}; our board evaluation follows the same principles.

\textit{Traffic analysis.} Sizes and timing of encrypted streams reveal motion~\cite{wampler2015} and which video is played~\cite{schuster2017}; traffic shaping hides activity at a large bandwidth cost~\cite{apthorpe2019}. No content encryption changes this, so we measure it separately.

We found no design that places a standardized AEAD inside a hardware JPEG encoder, keeps every level playable while authenticating each frame, hides block positions without a frame buffer, and closes the loop from the receiver back to the encoder.
